%
% reversi_220_persistenza_output.tex
%
% This file is part of the reversi program
% http://github.com/rcrr/reversi-v2
%
% Author Roberto Corradini mailto:rob_corradini@yahoo.it
% Copyright 2026 Roberto Corradini. All rights reserved.
%
% License
%
% This program is free software; you can redistribute it and/or modify it
% under the terms of the GNU General Public License as published by the
% Free Software Foundation; either version 2, or (at your option) any
% later version.
%
% This program is distributed in the hope that it will be useful,
% but WITHOUT ANY WARRANTY; without even the implied warranty of
% MERCHANTABILITY or FITNESS FOR A PARTICULAR PURPOSE. See the
% GNU General Public License for more details.
%
% You should have received a copy of the GNU General Public License
% along with this program; if not, write to the Free Software
% Foundation, Inc., 59 Temple Place - Suite 330, Boston, MA  02111-1307, USA
% or visit the site <http://www.gnu.org/licenses/>.
%

%-----------------------------------------------------------------------------
% Capitolo: Persistenza e output.
%
% Il dettaglio della fase finale della pipeline (i livelli SAVE,
% CLIENT_CHECK, PLOT e FINAL_LOG_PRINT): il file del modello, la verifica
% del round-trip attraverso il client, i grafici di diagnostica e il log
% consolidato della run.
%-----------------------------------------------------------------------------

\chapter{Persistenza e output}

Questa è l'ultima delle cinque fasi in cui si raggruppano i 19 livelli della
pipeline (tabella~\ref{tab:pipeline-levels} del capitolo 6): la fase
\emph{Persistenza e output}. I suoi quattro livelli — \texttt{SAVE},
\texttt{CLIENT\_CHECK}, \texttt{PLOT} e \texttt{FINAL\_LOG\_PRINT} — non
ricalcolano nulla: prendono ciò che la run ha già prodotto (il vettore di
pesi $w$, le serie di validazione, il registro degli eventi) e lo
trasformano negli output durevoli: il file del modello, la prova che il
modello è usabile da un client, i grafici di diagnostica e il log
consolidato.

\section{\texttt{SAVE}: il file del modello}

Il livello scrive un \emph{unico} file, sul percorso configurato
\texttt{base\_dir}/\texttt{output}. Il file è binario, a record
\emph{framed} (ogni valore è delimitato da due marcatori MAGIC
identici), little-endian e quindi leggibile invariabilmente da codice
C/Rust, e viene compresso con lz4-frame: la compressione è deterministica,
lo stesso contenuto produce sempre gli stessi byte. L'header è
\texttt{RGLM\_BINARY\_WEIGHTS} con versione 2: la v2 è
\emph{auto-descrivente}, perché incorpora la \emph{provenienza della
configurazione} — il record di configurazione normalizzato (lo stesso
stringa che il vecchio record laterale portava) più la sua impronta
sha3-256 — come prima sezione del corpo. Il file del modello è così
l'\emph{unica} fonte della configurazione della run: il record laterale
(\texttt{<output>.config.json}) è ritirato e non viene più generato.

Dopo la provenienza, il corpo porta, in ordine: i metadati (nome,
\texttt{empty\_count}, clipping del logit), le statistiche della
popolazione di addestramento (conto, media, deviazione standard,
varianza), l'informazione di ottimizzazione (motivo di stop, iterazioni,
valore finale, norma del gradiente), le metriche di validazione (campioni,
RMSE, MAE, loss), la definizione completa del feature set (intercetto,
mobility, pattern) e infine i dati dell'inferenza: la mappa degli offset di
feature \texttt{iwmap\_feature\_offset} e il vettore di pesi denso
\texttt{w\_dense}.

\begin{lstlisting}[language=python,
caption={L'ordine di scrittura del file (estratto di \texttt{model\_weights\_write\_file}).}]
with BinaryWriter(path, compressed=True) as w:
    w.write_header("RGLM_BINARY_WEIGHTS", version=2)
    # 1. provenienza della configurazione (sezione guida del corpo):
    w.write_string(ctx.config_checksum)
    w.write_string(ctx.cfg.model_dump_json(indent=4))
    # 2. metadati, 3. statistiche di addestramento,
    # 4. info di ottimizzazione, 5. metriche di validazione,
    # 6-8. feature set (intercetto, mobility, pattern),
    # 9. dati dell'inferenza:
    w.write_string(ctx.cfg.name)
    w.write_u32(ctx.cfg.empty_count)
    ...
    w.write_array(ctx.iwmap_feature_offset)
    w.write_array(ctx.w_dense)
\end{lstlisting}

Alla chiusura, il writer produce anche il sidecar \texttt{.SHA3-256}
(servizio di checksum attivo per default, nel formato standard
\texttt{sha3\_256sum}); il worker ne registra a livello \texttt{DEBUG} le
dimensioni e l'impronta — l'eco di integrità: ogni altro artefatto della
run ha già la propria riga di checksum nel log, il file finale è l'unico a
riceverla qui. Il passo \texttt{down} è un no-op: il worker è stateless.

\section{\texttt{CLIENT\_CHECK}: la verifica del round-trip}

Il modello salvato non è utile finché non è provato \emph{usabile} da un
client: questo è il lavoro di \texttt{CLIENT\_CHECK}. Il livello rilegge il
file appena scritto con \texttt{model\_weights\_read\_file} (il lettore v2
ri-verifica per strada l'impronta della configurazione incorporata) e
ottiene il dataclass \texttt{ModelWeights}: i pesi, la mappa degli offset,
il feature set — tutto ciò che a un client serve.

Costruisce poi la \emph{Evaluation Function} vettorizzata del client,
\texttt{make\_evaluation\_function\_array}: la funzione è chiusa su due
artefatti precalcolati — gli offset di avvio di \texttt{w\_dense} per
istanza (l'array dei confini di feature espanso alla granularità delle
istanze) e la mappa affine $\mathrm{z}\to\mathrm{y}$ — e valuta l'intero
set di validazione in una chiamata: indici principali, raccolta dei pesi
densi, somma nel predittore lineare, sigmoid, ritorno nello spazio Y.

Il cuore del controllo è un confronto di residui: i residui di validazione
ricalcolati dai pesi del \emph{file}, in spazio Y (vista f64 senza perdita
per arrotondamento), e la serie di riferimento \texttt{vld\_rn} pubblicata
dal livello \texttt{VALIDATE} — nessun artefatto di validazione muta tra i
livelli 14 e 17, dunque l'array in memoria \emph{è} il riferimento. La
massima differenza deve restare sotto la tolleranza
\texttt{client\_check\_tol}; un superamento fa fallire la run con un
\texttt{RuntimeError}, loggato a livello \texttt{ERROR} con l'entità della
divergenza e la prima posizione in cui compare.

\begin{lstlisting}[language=python,
caption={Il confronto del round-trip (estratto di \texttt{wrk\_client\_check}).}]
mw = model_weights_read_file(output_path, compressed=True)
ef = make_evaluation_function_array(mw)
y_pred_file = ef(ctx.vld_positions)
rn_file = (y_pred_file.astype(np.float64)
           - ctx.vld_game_values.astype(np.float64))
rn_ref = ctx.vld_rn
max_diff = float(np.abs(rn_file - rn_ref).max())
if max_diff >= ctx.cfg.client_check_tol:
    raise RuntimeError("CLIENT_CHECK failed")
\end{lstlisting}

In sintesi: modello in memoria $\to$ file $\to$ client $\to$ gli stessi
numeri. Se il round-trip torna identico entro la tolleranza, il file è
dichiarato affidabile.

\section{\texttt{PLOT}: i grafici di diagnostica}

Il livello produce le figure \emph{matplotlib} della run. Un \emph{pilot} a
due stadi li governa: l'interruttore \texttt{plot.enable} (se false, il
livello è un no-op pulito anche con nomi in lista) e la \emph{whitelist}
\texttt{plot.plots}, che sceglie quali dei sette nomi canonici generare: un
nome in lista è scritto, uno assente è saltato, una lista vuota rende il
livello inerte.

I grafici non rimisurano nulla: \texttt{hlp\_plot\_collect\_data} assembla
le cinque serie f64 che la pipeline ha già calcolato — i valori di gioco e
i residui di addestramento (il passo forward denso di
\texttt{ANALYTICS}) e i valori, le stime e i residui di validazione (di
\texttt{VALIDATE}) — tutte in vista f64 senza perdita, così che nessun
rumore di arrotondamento f32 entri nelle figure. La directory di
destinazione è \texttt{plot.dir\_name}: usata come sta se assoluta (può
vivere fuori da \texttt{base\_dir}), risolta su \texttt{base\_dir} se
relativa (default: \texttt{plots}). Ciascun grafico è salvato in ogni
formato della lista \texttt{plot.format} (la coppia \texttt{"pdf,png"}
scrive entrambi: il pdf vettoriale e il png alla risoluzione
\texttt{plot.dpi}).

L'aspetto e la \emph{determinismo} sono guidati da tredici manopole (il
dataclass \texttt{HlpPlotStyle}, che la sezione \texttt{plot} della
configurazione porta con sé): i semi per il jitter e il
sotto-campionamento del KDE, le grandezze del jitter, i domini degli assi
(gli assi di valore di gioco sono sempre il dominio completo $[-64,
+64]$), il passo dei tick e il posizionamento delle etichette delle serie.
Ogni casualità è guidata da generatori con seme fisso, quindi le figure
sono riproducibili di run in run.

\begin{table}[htbp]
\centering
\small
\begin{tabularx}{\textwidth}{lX}
\toprule
Nome & Contenuto \\
\midrule
\texttt{residual\_density} & Densità (KDE) di tre serie di residui di valore di gioco: di riferimento (i valori di addestramento centrati sulla media), di addestramento e di validazione; la riga tratteggiata segna la media dei residui di addestramento. \\
\texttt{resid\_vs\_fitted} & Nuvola dei residui di validazione contro i valori stimati, con jitter a seme fisso su entrambi gli assi. \\
\texttt{resid\_vs\_gv} & Nuvola dei residui di validazione contro i valori di gioco effettivi. \\
\texttt{gv\_by\_egv\_box} & Box plot dei valori di gioco di validazione, raggruppati per valore stimato arrotondato alla griglia a passo 2. \\
\texttt{resid\_mean\_by\_egv} & Media del residuo con barre $\pm$ deviazione standard per valore stimato, contro la linea a zero: la vista dei bias sistematici per livello. \\
\texttt{parity} & Parità: valori di gioco di validazione contro valori stimati, con la retta a 45$^\circ$; la qualità del fit globale a colpo d'occhio. \\
\texttt{resid\_qq} & QQ plot normale dei residui di validazione: quantili campionario contro quantili teorici standard-normali, con la riga di scala $y=\mathrm{sd}\,x$. \\
\bottomrule
\end{tabularx}
\caption{I sette grafici canonici della sezione \texttt{plot}.}
\label{tab:plot-names}
\end{table}

\section{\texttt{FINAL\_LOG\_PRINT}: il log consolidato}

Ultimo livello, ultimo output: il log della run. Se
\texttt{log\_file\_name} è assente nella configurazione, il livello è un
no-op pulito (il file non viene generato); altrimenti
\texttt{hlp\_final\_log\_print\_write\_log} renderizza il \emph{registro
eventi della macchina} — la lista di \texttt{LogEntry} che la
\texttt{StateMachine} riempie a ogni \texttt{log\_event},
indipendentemente dalla gate di verbosità della console — a partire da
\texttt{DEBUG} (il dettaglio delle hot loop a livello \texttt{TRACE} resta
fuori dal file) e lo scrive in UTF-8 su \texttt{base\_dir}/\texttt{log\_file\_name},
creando la directory di destinazione se serve. Il file è il log completo
fino al momento della scrittura: contiene la riga ``Starting worker
\texttt{FINAL\_LOG\_PRINT}'', ma non la ``Completed UP step'' (loggata
dopo il ritorno). Il passo \texttt{down} è un no-op: il file non viene
rimosso.

Con i quattro livelli la fase — e la pipeline — si chiude: la run è
completa, dal file di configurazione al modello salvato, verificato,
tracciato in grafici e registrato nel log, pronto a essere caricato e
usato.
